\chapter{Introduction}\label{ch:introduction}

\section{Context}
Deep neural network inference is increasingly embedded in software systems that must satisfy not only predictive requirements, but also operational ones such as latency, deployability, maintainability, scalability, and isolation. Edge-intelligence and split-computing surveys describe this shift as part of a broader move from purely centralized inference toward device, edge, and cloud collaboration~\cite{edge_intelligence,split_computing_survey}. In many practical deployments, inference still runs as a monolithic component inside a single process. That architecture minimizes boundary-crossing overhead, but it also couples model execution into one deployment unit and limits architectural flexibility. A microservice-oriented design offers a different trade-off: the inference pipeline can be decomposed across service boundaries, allowing clearer separation of responsibilities, independent deployment, and more explicit control over resource placement.

For deep learning systems, however, decomposition is not only a software architecture decision. It is also a systems problem shaped by the structure of the neural network itself and by the environment in which the resulting services run. When a model is split across an execution boundary, intermediate activations must be serialized, transferred, and reconstructed before downstream computation can continue. Prior work on split computing and distributed DNN inference shows that these costs are highly sensitive to where the boundary is placed. DNNSplit demonstrates that split points should be selected systematically with respect to latency and cost rather than by arbitrary layer choice \cite{dnnsplit}. DeeperThings and survey work on DNN partitioning likewise show that early layers often produce large feature maps that are expensive to transmit, while later boundaries may reduce transfer volume but shift the compute balance too far toward one side of the partition \cite{deeperthings, survey_partitioning}.

The deployment environment introduces an additional dimension to this problem. A split that is acceptable in a controlled local benchmark may behave differently once it is moved into containerized, orchestrated, or cloud-hosted execution, because container platforms, managed Kubernetes services, and cloud infrastructure add their own scheduling, networking, and virtualization effects~\cite{managed_k8s_perf,virtualization_costs,cni_k8s_ic2e21}. Beyond performance, there is also a security dimension. As inference services move into shared or remote infrastructure, communication-level protections such as service identity and mutual TLS become relevant because they strengthen inter-service trust boundaries, while service-mesh studies show that this protection can add measurable runtime and resource cost~\cite{meshinsight_socc23,bremlerbarr2025mtls}. Mutual TLS terminates at the pod boundary, however, so intermediate activations are still exposed to the runtime environment of the downstream service; prior work shows that an untrusted remote partition holding only those activations can reconstruct the input~\cite{he2021_attacking_collaborative}, motivating an additional layer of protection. Confidential-computing techniques such as trusted execution environments may add such a layer, but their overhead and protection boundary are workload- and platform-dependent~\cite{guanciale2022_confidential_quartet,mo2024_ml_confidential_computing}, and recent work shows that confidential VMs still have residual attack surfaces beyond memory encryption, including malicious-host interrupt injection~\cite{schluter2024_heckler}. This makes microservice-based inference a broader software engineering problem involving architectural partitioning, deployment infrastructure, and security mechanisms rather than only local latency optimisation.

\section{Problem Statement}
The central problem addressed in this thesis is how to design and evaluate deep neural network inference as a microservice system without losing control over performance. A poor split can make a distributed design clearly inferior to monolithic execution. Prior partitioning systems show that if the boundary is placed too early, the system may incur excessive latency because large activation tensors must cross the service interface; if the boundary is placed too late, communication cost decreases, but the downstream service may be left with too little computation to justify the split~\cite{neurosurgeon,deeperthings,dnnsplit}. In other words, minimizing activation size alone is not enough; the chosen boundary must also preserve a meaningful compute distribution across the resulting services.

That architectural problem is only the first layer of the thesis. After a defensible split region has been identified locally, the same decomposition logic must also be evaluated under progressively more realistic deployment conditions. Container orchestration and cloud deployment introduce additional scheduling, networking, and platform overhead that may alter the trade-offs seen in a local benchmark~\cite{managed_k8s_perf,noise_in_clouds}. Finally, if the selected deployment path is security-hardened, the system may incur further overhead in exchange for stronger inter-service authentication, encrypted communication, and confidential execution for selected services~\cite{meshinsight_socc23,yan2023_cvm_overheads}. The thesis therefore treats the problem as a staged evaluation pipeline: first identify and explain suitable boundaries under controlled local conditions, then study how a coarse-stage service-chain family behaves in Kubernetes and Azure deployments, and finally measure the impact of adding communication-level and VM-level security hardening to the selected Azure deployment path.

\section{Research Questions}
This thesis is organised around two broader research strands. RQ1 studies how architectural split choice and deployment context affect microservice-based inference. RQ2 studies the performance, operational, and protection-scope trade-offs introduced when the selected Azure deployment path is security-hardened. The research questions are:

\begin{description}
    \item[\textbf{RQ1.1}] Which coarse architectural stage boundaries provide the most suitable two-part microservice decompositions of ResNet-18 under controlled local execution, and how do their latency and boundary-crossing costs compare to monolithic inference?
    \item[\textbf{RQ1.2}] How does finer-grained refinement within the late-stage coarse region carried forward from RQ1.1 affect latency and communication-related overhead for two-part microservice inference?
    \item[\textbf{RQ1.3}] How do activation-transfer size and compute distribution explain the latency and boundary-crossing behaviour observed for the decomposition choices evaluated in RQ1.1 and RQ1.2?
    \item[\textbf{RQ1.4}] How does increasing the number of microservices affect end-to-end inference cost when ResNet-18 is decomposed into synchronously chained coarse architectural stages under in-cluster Kubernetes execution?
    \item[\textbf{RQ1.5}] Does the local Kubernetes latency pattern observed for the coarse-stage ResNet-18 microservice chain transfer to Azure Kubernetes Service, and how do the absolute and relative costs change under managed cloud execution?
    \item[\textbf{RQ1.5b}] When every chain hop is forced across a node boundary, how much additional cost is introduced relative to the same-node RQ1.5 baseline?
    \item[\textbf{RQ2.1}] What performance and operational overhead is introduced when the selected AKS microservice inference path is hardened with service identity, mutual TLS, and inter-service authorization policy?
    \item[\textbf{RQ2.2}] What additional performance and operational overhead is introduced when VM-level confidential execution is added to the already communication-secured AKS inference chain, when the downstream protected service is placed on AMD SEV-SNP confidential nodes?
    \item[\textbf{RQ2.3}] Which security-hardening configuration provides the most appropriate trade-off between performance cost, operational complexity, and protection scope for the selected AKS microservice inference deployment?
\end{description}

RQ1.1 and RQ1.2 form the local split-selection stages. RQ1.3 is an explanatory synthesis that uses activation-transfer size and internal compute distribution to interpret those local results. RQ1.4 and RQ1.5 then evaluate the same coarse-stage chain family in local Kubernetes and AKS, while RQ1.5b acts as a sensitivity stage for cross-node placement. RQ2.1 adds communication-level hardening through service identity, mutual TLS, and authorization policy. RQ2.2 adds selective AMD SEV-SNP VM-level confidential execution for the downstream protected service. RQ2.3 synthesises those security stages against the plain AKS baseline.

\section{Approach}
The thesis uses a staged experimental methodology centred on ResNet-18, a residual architecture whose stage structure provides natural coarse split candidates~\cite{resnet}. The initial benchmark compares one monolithic baseline against four two-part split configurations placed after \texttt{layer1}, \texttt{layer2}, \texttt{layer3}, and \texttt{layer4}. These boundaries correspond to major architectural stages in the network and therefore provide a principled coarse sweep from early to late splits. This first stage is executed locally with CPU-only inference, FP32 precision, gRPC over localhost, fixed warmup and measurement windows, and functional parity checks between monolithic and split execution so that boundary effects can be isolated under controlled conditions, following systems-benchmarking guidance that emphasises controlled comparisons and explicit reporting of measurement context~\cite{benchmarking_sc15}.

The local selection stages are then refined and explained before the deployment context is changed. A block-level split inside the carried-forward region tests whether finer granularity changes the conclusion, while an internal timing study of the monolithic model supports the interpretation of activation-transfer and compute-distribution effects. The same model and benchmark contract are then carried into a predefined Kubernetes service-chain family, from a one-service monolithic deployment to a five-service coarse-stage chain, and this family is evaluated first in local Kubernetes and then on Azure Kubernetes Service.

The final stages add security hardening to the selected Azure path. RQ2.1 compares the plain AKS chain against a managed-Istio configuration with service identity, mutual TLS, and inter-service authorization policy. RQ2.2 keeps that communication-security contract fixed and moves the downstream protected service onto AMD SEV-SNP confidential nodes to measure the incremental cost of VM-level confidential execution. RQ2.3 then compares the plain, communication-secured, and selectively confidential configurations. This design follows the broader literature showing that DNN partitioning must be evaluated as a trade-off between communication and computation rather than as a purely architectural choice~\cite{ga_partitioning, fine_grained_partitioning, delay_aware_inference}. In this thesis, that trade-off is examined across multiple deployment and protection layers rather than being treated as a purely local benchmarking problem.

\section{Outline}
The remainder of this thesis is organised as follows. Chapter~\ref{ch:background} introduces the technical background required for understanding deep neural network inference, ResNet-style architectures, microservice-based deployment, and the broader systems context in which distributed inference is evaluated. Chapter~\ref{ch:related} reviews prior literature on split computing, DNN partitioning, distributed inference, and the trade-offs that motivate this thesis. Chapter~\ref{ch:methods} describes the research methodology and the staged evaluation design used to move from local boundary screening to broader deployment and security contexts. Chapter~\ref{ch:experiments} presents the empirical evaluation stages, beginning with RQ1.1 and continuing through the later deployment-oriented and security-hardening analyses. Chapter~\ref{ch:analysis} interprets those findings in relation to the research questions and discusses their implications and limitations. Finally, Chapter~\ref{ch:conclusion} summarises the thesis and outlines directions for future work.
